% !TEX root = ../main.tex
\section{Conclusiones}
\label{sec:conclusiones}

El análisis matemático, el desarrollo declarativo en \texttt{MiniZinc} de los modelos de Parte 1 y Parte 2 adherido rigurosamente a los fundamentos del curso, la construcción de la interfaz gráfica y normalizada de entrada/salida y la experimentación exhaustiva sobre trece instancias permiten establecer las siguientes conclusiones técnicas debidamente sustentadas:

\begin{enumerate}
    \item \textbf{Idoneidad de la canalización primal-dual sin salir del núcleo declarativo del curso:}
    Representar simultáneamente el cronograma mediante el arreglo por turnos $Orden_k$ y su permutación inversa por escenas $Pos_e$ a través del predicado global \texttt{inverse}, acoplados con un vector de tiempos acumulados $Acum_k$, permite expresar el ingreso, el retiro y la estadía continua de cada actor empleando únicamente agregadores \texttt{min}, \texttt{max} y restricciones \texttt{element} sobre conjuntos estáticos de escenas propias $E_a$, mientras todo el preprocesamiento de rangos salariales y duraciones ordenadas se resuelve con listas por comprensión y sumas aritméticas elementales. Esta arquitectura evita tanto la multiplicación cuadrática de variables como la reificación superflua sobre las restricciones primarias ($C_1$--$C_5$).

    \item \textbf{Exactitud y eficiencia de la escalarización lexicográfica en optimización jerárquica:}
    En la Parte 2, el requerimiento de minimizar como criterio secundario la coincidencia en el estudio de los actores incluidos en \texttt{Evitar} sin deteriorar bajo ninguna circunstancia el costo salarial mínimo se resuelve de forma exacta en una única ejecución mediante la función lineal $ObjetivoLex = (\overline{TC} + 1)\cdot Costo + TiempoCompartido$. La contrastación contra el oráculo independiente confirma que esta formulación garantiza prioridad lexicográfica absoluta y mantiene dominios enteros reducidos libres de desbordamiento numérico.

    \item \textbf{Impacto decisivo y complementario de la eliminación de simetrías y las restricciones redundantes:}
    La comparación experimental directa frente al modelo primario original ($M_0$) demuestra que romper la inversión temporal sobre las escenas activas ($S_1$), ordenar escenas gemelas ($S_2$) y fijar en prefijo las escenas sin actores ($S_3$) preserva intacto el valor óptimo (permitiendo reconstruir algebraicamente toda la órbita de soluciones equivalentes descartadas) mientras poda más del $51\%$ de los nodos explorados y fallidos en instancias asimétricas generales y hasta un $94{,}1\%$ ($17{,}04\times$ por simetrías solas y $24{,}61\times$ en conjunto) cuando coexisten escenas vacías y gemelas (\texttt{Instancia01\_P1}). A su vez, las restricciones redundantes de intermediación lineal con coeficientes constantes ($R_1$), diferencias temporales entre subintervalos ($R_2$), amplitud posicional ($R_3$) y exceso mínimo ($R_4$) reducen conjuntamente los nodos explorados y fallidos en $10{,}13\times$ (\texttt{Instancia03\_P1}) y $11{,}45\times$ (\texttt{Instancia03\_P2}), haciendo posible certificar la optimalidad global de ``Desenfreno de Pasiones'' ($M=20$, $2{,}43 \times 10^{18}$ permutaciones) tanto en la Parte 1 ($Costo^* = 871$ en $6{,}60\text{ s}$) como en la Parte 2 ($Costo^* = 880, TC^* = 37$ en $1{,}57\text{ s}$).

    \item \textbf{Necesidad crítica de alinear la estrategia de distribución con la topología del acoplamiento temporal:}
    El análisis estadístico sobre nodos explorados, nodos fallidos y tiempos de respuesta de las ocho estrategias ($E_1$--$E_8$) evidencia que las heurísticas genéricas (cronológica izquierda-a-derecha $E_3$, \emph{first-fail} $E_4$--$E_5$ o la búsqueda libre por defecto del motor $E_8$, la cual queda estancada en costos altamente subóptimos de $1046$ en Parte 1 y $1094$ en Parte 2) sufren una degradación severa al crecer $M$, dado que dejan abiertos los extremos de permanencia actoral hasta niveles profundos del árbol. En contraste, instanciar las posiciones desde ambos extremos hacia el centro (\emph{Outside-In}) probando primero las escenas de menor peso salarial ($E_1$ y $E_7$) activa la ruptura de simetría $S_1$ en profundidad $2$, cierra tempranamente los intervalos actorales y descubre incumbentes casi óptimos desde las primeras ramas, consolidándose junto con el aprendizaje de cláusulas de \texttt{Chuffed} como la estrategia superior en toda la batería evaluada.
\end{enumerate}
